% Placeholder LaTeX skeleton for the deep research paper.
% Choose the final template (acmart / IEEEtran / usenix) after venue is decided.
% Current draft targets a venue-agnostic 12-page article layout.

\documentclass[11pt]{article}
\usepackage[margin=1in]{geometry}
\usepackage{amsmath,amssymb}
\usepackage{booktabs}
\usepackage{graphicx}
\usepackage{hyperref}
\hypersetup{
  pdftitle={Golden Signals for AI Workloads on Shared GPU Clusters},
  pdfauthor={Andrew Espira and Sharath Kumar},
  colorlinks=true,linkcolor=blue,urlcolor=blue,citecolor=blue
}

\title{Beyond the Four Golden Signals:\\
Observability Primitives for AI Workloads on Shared GPU Clusters}
\author{Andrew Espira \and Sharath Kumar\\
Department of Data Science, Saint Peter's University}
\date{Draft --- \today}

\begin{document}
\maketitle

\begin{abstract}
The four golden signals of monitoring --- latency, traffic, errors, and
saturation --- were codified for stateless request/response services.
When AI training and inference are placed on shared GPU clusters, each of
the four either loses precision or actively misleads: GPU utilization is a
well-documented lying signal, HTTP-style errors miss silent quality
regressions, request latency does not capture queue wait or token-level
tails, and QPS is undefined for training. We propose five golden signals
tailored to AI GPU workloads --- \emph{Queue-of-Work Reliability},
\emph{GPU-Compute Effectiveness}, \emph{HBM-Pressure Headroom},
\emph{Straggler Amplification Index}, and \emph{Output-Quality
Calibration} --- each defined precisely, each with an SLO template and a
burn-rate multi-window alert, and each validated on measured data from
four heterogeneous environments (Amazon EKS, AWS ParallelCluster Slurm,
Alibaba GPU trace v2020/2023, Google Borg 2019). We show the five signals
are mutually informative rather than redundant and that they are cheaply
computable from telemetry a well-instrumented cluster already emits.
\end{abstract}

\section{Introduction}\label{sec:intro}
\emph{TODO. Motivating example: an EKS cluster with 100\% \texttt{util\%}
and healthy p99 HTTP latency is silently dropping 22\% goodput to
stragglers. Contribution list.}

\section{Background}\label{sec:bg}
\emph{TODO. Golden signals lineage (SRE Book, SRE Workbook, Implementing
SLOs). Assumptions of the classical four and where they break for GPU AI
workloads (\S~\ref{sec:critique}).}

\section{Why the classical four travel poorly}\label{sec:critique}
\emph{TODO. See \texttt{docs/CRITIQUE.md} for the full argument to be
lifted into this section.}

\section{Five golden signals, defined}\label{sec:signals}
\emph{TODO. See \texttt{docs/SIGNALS\_TAXONOMY.md} for the taxonomy.
Each subsection covers one signal.}

\subsection{Queue-of-Work Reliability (QWR)}\label{sec:qwr}
\subsection{GPU-Compute Effectiveness (GCE)}\label{sec:gce}
\subsection{HBM-Pressure Headroom (HPH)}\label{sec:hph}
\subsection{Straggler Amplification Index (SAI)}\label{sec:sai}
\subsection{Output-Quality Calibration (OQC)}\label{sec:oqc}

\section{Data and measurement protocol}\label{sec:data}
\emph{TODO. Four environments; per-signal extraction recipe cross-linked
to \texttt{code/}.}

\section{Empirical validation}\label{sec:results}
\emph{TODO. Per-signal behavior across four environments;
cross-signal correlation table; failure-mode case studies;
alert-simulation comparison to classical golden-signal alerts.}

\section{Operational integration}\label{sec:ops}
\emph{TODO. SLO templates, tier-qualification tie-in, calibration monitor
as a shared sub-component of three of the five signals.}

\section{Discussion and Limitations}\label{sec:disc}
\emph{TODO.}

\section{Related Work}\label{sec:rw}
\emph{TODO. See \texttt{docs/RELATED\_WORK.md}.}

\section{Conclusion}\label{sec:conc}
\emph{TODO.}

\section*{Use of AI Assistance}
Draft text of this paper was edited with the help of a large language model
(Anthropic Claude). All results, tables, figures, methodological choices,
and interpretations were authored and verified by the human authors; the
LLM did not generate experimental results.

\section*{Acknowledgments}
The framing of this work --- that AI workloads on shared GPU clusters
lack a widely-accepted golden-signal set, and that codifying one is a
first-class SRE problem rather than a monitoring-tool problem ---
draws on the Google SRE community's public discussions on monitoring
AI workloads in large production deployments. We thank the Saint
Peter's University Department of Data Science for infrastructure and
guidance.

\end{document}
